\documentclass[11pt,a4paper]{article}
\usepackage[margin=1in]{geometry}
\usepackage{amsmath,amssymb,amsthm}
\usepackage{algorithm}
\usepackage{algpseudocode}
\usepackage{booktabs}
\usepackage{hyperref}

\theoremstyle{definition}
\newtheorem{definition}{Definition}[section]
\newtheorem{theorem}{Theorem}[section]
\newtheorem{lemma}[theorem]{Lemma}
\newtheorem{remark}{Remark}[section]

\title{\textbf{Rigorous Mathematical Specification, Boltzmann Distribution, and Entropy Limits of Temperature Sampling}}
\author{Team Scaibu \\ \small Email: \texttt{jaydeep.9664920749@gmail.com} \\ \small Architecture and Core Engine Team}
\date{October 2026}

\begin{document}
\maketitle

\begin{abstract}
Autoregressive decoding converts unnormalized logit vectors into token predictions. Temperature sampling controls the sharpness of the categorical probability distribution via annealing parameter $T \ge 0$. As $T \to 0$, the Boltzmann distribution converges to the Dirac delta distribution at the maximum logit (greedy decoding); as $T \to \infty$, the distribution converges to the discrete uniform distribution. This specification formalizes the thermodynamic temperature mapping, proves asymptotic limits, and details numerical verification.
\end{abstract}

\section{Mathematical Formulation}

Let logits $\mathbf{z} \in \mathbb{R}^V$ and temperature $T \ge 0$.
\begin{enumerate}
    \item If $T = 0$ (or $T < 10^{-6}$):
    \begin{equation}
    P(v) = \begin{cases} 1.0, & \text{if } v = \arg\max_{j} z_j \\ 0.0, & \text{otherwise} \end{cases}
    \end{equation}
    \item If $T > 0$:
    \begin{equation}
    P(v; T) = \frac{\exp((z_v - \max_j z_j) / T)}{\sum_{j=0}^{V-1} \exp((z_j - \max_l z_l) / T)}
    \end{equation}
\end{enumerate}

\section{Theoretical Guarantees: Asymptotic Thermodynamic Limits}

\begin{theorem}[Zero and Infinite Temperature Limits]
For distinct maximum logit $z^* > z_j$ ($\forall j \ne *$):
\begin{align}
\lim_{T \to 0^+} P(v; T) &= \delta_{v, *} \\
\lim_{T \to \infty} P(v; T) &= \frac{1}{V} \quad \forall v \in \{0, \dots, V-1\}
\end{align}
\end{theorem}

\begin{thebibliography}{9}
\bibitem{ackley1985} D.~H.~Ackley, G.~E.~Hinton, and T.~J.~Sejnowski, ``A learning algorithm for Boltzmann machines,'' \emph{Cognitive Science}, vol.~9, no.~1, pp.~147--169, 1985.
\end{thebibliography}

\end{document}
